\subsection{Evaluation}\label{sec:2.3}

% Sources: DETECTION_TASK.md; INTENDED_USE.md; UCSF_PLAN.md; SYSTEM_PLAN.md; METHODS_PLAN.md; NUMBERS.md 2.1;
% prism/analyses/detection.py (fold_ids), prism/analyses/system.py (locked_models, _fit); refs varoquaux2018.
\paragraph*{Development and testing without retraining.} Each cohort had one role: MU-Glioma-Post for development,
LUMIERE and UCSF-ALPTDG for external testing (Table~\ref{tab:eval}). On MU we evaluated every model by five repeats of
fivefold cross validation grouped by patient. In each repeat the patients were split into five folds, a model was
fitted on the scans of four folds and scored the scans of the fifth, so that every scan was scored once by a model that
had not seen its patient. Models were fitted on all grade 4 detection scans of the training folds, and the AUC was
computed over the held-out scores of the 146 scans in the intended use and averaged over the five repeats. We repeated
the split because a single split of a small cohort gives an estimate with wide error \citep{varoquaux2018}. We then
fitted each model once on all 181 grade 4 detection scans of MU, including the 35 without a documented start of
radiotherapy by the scan day, and applied it to LUMIERE and UCSF-ALPTDG without retraining or recalibration. At
UCSF-ALPTDG every model was fixed in a written plan before any score was compared with the label. At LUMIERE the
combined model and its comparison with the time-and-enhancement model were fixed in advance, and the other scores are
exploratory.

% Sources: CLINICIAN_BASELINE_PLAN.md; METHODS_RESULTS.md E4; prism/analyses/clinician_baseline.py;
% prism/analyses/timing_only.py and docs/results/TIMING_ONLY.md (timing alone); ref beaulieujones2021.
\paragraph*{The time-and-enhancement model.} We compared every score with a reference model built from what a
clinician reading the scan already knows: when the scan was taken relative to radiotherapy, how much enhancing tumour
it shows, and whether that tumour has grown. Models can reach a high AUC by reproducing what clinicians already know
\citep{beaulieujones2021}, so the difference in AUC between a score and this model measures what the score adds. The
time-and-enhancement model is a ridge logistic regression on three inputs: the time window after radiotherapy, the log
enhancing volume, and its log ratio to the previous scan. On MU the window enters as indicators for before
radiotherapy, during it and within 12 weeks after its end, with later scans as the reference; the model is fitted in
the same folds and with the same fitting procedure as PRISM's models. Externally it takes the 12-week flag in place of
the window and a flag for whether a previous scan exists, and is fitted once on MU and applied unchanged like every
other model. At LUMIERE the flag is the rater's mark for scans within three months of radiotherapy. At UCSF-ALPTDG the
window had to be approximated (below), so there the model was applied only in the segmentation comparison, outside the
planned comparison. The reading was fixed in advance: a gain whose paired 95\% CI lies above zero counts as a gain,
and one whose CI includes zero is reported as no clear gain. To show how far time alone ranks scans, we also scored
them by timing alone (exploratory, Table~\ref{tab:models}): the days since the end of radiotherapy and the visit
number, each used as a score without a fitted model, and a ridge model on these two, the days since diagnosis and the
12-week flag. Externally, where these dates are missing, the unfitted score is the days since surgery (LUMIERE) or
since the start of radiotherapy (UCSF-ALPTDG). A difference in AUC does not show which pairs of scans a score ranks
well, so we also split each AUC by time window.

% Sources: METHODS_PLAN.md M1; METHODS_RESULTS.md M1 and UCSF section; prism/analyses/methods.py (split_auc,
% m1_stats, boot); prism/analyses/ucsf_external.py (rt_phase); refs kallus2019, janes2008, janes2009.
\paragraph*{Splitting the AUC by time window.} The AUC is the share of pairs, one scan with progression and one
without, that a score ranks in the right order. We split these pairs into two groups. In a pair within a window, both
scans were taken in the same window after radiotherapy, for example both later than 12 weeks after its end; the window
cannot tell them apart, so the score must rank them on other evidence. In a pair across windows, the scans come from
different windows, and the window alone can often order them, because the share of scans with progression differs
between windows. The split shows the reader how much of a score's AUC comes from each kind of pair. Let $P$ and $N$ be
the scans with and without progression, $s$ a score, and $c$ the time window of each scan. The AUC is
\begin{equation}
\mathrm{AUC} = \frac{1}{|P|\,|N|} \sum_{i \in P} \sum_{j \in N} K(s_i, s_j), \qquad
K(a, b) = \mathbb{1}[a > b] + \tfrac12\, \mathbb{1}[a = b].
\end{equation}
The pairs split into $W = \{(i, j) : c_i = c_j\}$, the pairs within a window, and its complement $B$, the pairs across
windows. With $A_W$ and $A_B$ the same average taken over $W$ and over $B$, and $w = |W| / (|P|\,|N|)$ the share of
pairs within a window,
\begin{equation}
\mathrm{AUC} = w\, A_W + (1 - w)\, A_B
\end{equation}
holds exactly; it is the decomposition of Kallus and Zhou \citep{kallus2019} with the pairs grouped into these two
sets. $A_W$ is the ranking with the window held fixed, a pair-weighted analogue of the covariate-adjusted AUC
\citep{janes2008,janes2009}. A score whose $A_B$ exceeds its $A_W$ ranks pairs across windows better than pairs within
one. To show how much of $A_B$ the window alone could supply, we also scored every scan by the share of scans with
progression in its window in the training data (the window-only score). Its $A_W$ is close to 0.5, and its $A_B$ is
the ranking the window gives by itself. At MU the window comes from the treatment dates: during radiotherapy, within 12
weeks after its end, and later. At LUMIERE it is the rater's flag for scans within three months of radiotherapy. At
UCSF-ALPTDG, which records only the start of radiotherapy, it assumes a six-week course, so that 42 to 125 days after
the start counts as within 12 weeks of the end. Externally the window-only score takes MU's share for the matching
window, and MU's overall share for scans without a dated course of radiotherapy. Scans without a value of a score are
left out of that score's pairs, and the identity was checked on every bootstrap sample.
